\vspace{-0.5em}
\section{Evaluation and Results}\label{sec:evaluation}

Our evaluation shows that PaTAS produces trust estimates that converge during training, respect the properties of \cref{sec:ptastheory}, and remain interpretable under noisy features, corrupted labels, and adversarial perturbations.

We conduct four experiments of increasing complexity: (1)~\emph{Breast Cancer Classification} on a small tabular medical dataset; (2)~\emph{MNIST Digit Classification} across architectures under controlled uncertainty; (3)~\emph{Poisoned MNIST} under adversarial corruption; and (4)~\emph{Comparative Trust Evaluation} against established uncertainty baselines on out-of-distribution and corruption detection, characterizing when input-level trust is informative.

\vspace{-0.5em}
\subsection{Experimental Setup }
\subsubsection{Experiment 1 - Breast Cancer Classification}\label{exp:1}
We use the \href{https://archive.ics.uci.edu/dataset/17/breast+cancer+wisconsin+diagnostic}{Breast Cancer Wisconsin (Diagnostic) Dataset}~\cite{breastcancer-wisconsin}, 569 samples with 30 numeric features (e.g., radius, area, symmetry) derived from breast mass fine needle aspirates, classified into benign or malignant. The network has 30 input, 16 hidden (ReLU), and 2 output (Softmax) neurons, trained for 15 epochs with batch size 64 and learning rate 0.2, reaching 97\% test accuracy on unmodified data.

We degrade the training data in controlled ways and assign corresponding trust assessments, combining three extreme profiles for features and labels alike, fully trusted \((1,0,0)\), fully distrusted \((0,1,0)\), and fully uncertain \((0,0,1)\), into nine cases; these dogmatic and vacuous opinions are the canonical SL boundary cases.

Degradation is simulated with controlled perturbation functions (Appendix~\ref{app:perturb}): uncertain features receive additive uniform noise clipped to the valid range, distrusted features are resampled uniformly, and labels are degraded analogously through random noise (uncertainty) or complete replacement (distrust).

Finally, two intermediate scenarios complement the boundary cases (last two rows of \cref{tab:cancer_results}): the fully uncertain features-and-labels setting re-assessed as the mixed opinion \((0.25,0.25,0.5)\), and a mild degradation where features are perturbed with probability 0.15 and assessed as \((0.25,0,0.75)\).

\subsubsection{Experiment 2 - MNIST}\label{exp:2}

We evaluate PaTAS on MNIST~\cite{lecun1998mnist} (60{,}000 training and 10{,}000 test \(28\times28\) images, ten classes) with four architectures of increasing width, 784-16-10 through 784-128-10; not state of the art, but sufficient to expose how trust propagates across model sizes. Each uses ReLU and Softmax, trained for 20 epochs with batch size 128 and learning rate 0.05, reaching 99\% test accuracy on fully trusted data; the gradient-evidence threshold is $\epsilon = 0.05$ here and in Experiments 3 and 4. Feature and label trust assessments are fully uncertain, corresponding to no knowledge of the dataset; for contrast we repeat the smallest architecture with a fully trusted assessment.

\subsubsection{Experiment 3 - Poisoned MNIST}\label{exp:3}
We poison one third of the MNIST training images by flipping the labels of digits 6 and 9 while adding a visible fixed-size patch at the top-left corner, following common backdoor practice~\cite{chen2017targetedbackdoorattacksdeep}; the remaining two thirds stay clean. Using the 784-128-10 architecture, patch pixels and the labels of patched 6s and 9s are assigned distrust and everything else trust, so PaTAS focuses on the corrupted regions while trusting the unaffected data.


\subsubsection{Experiment 4 - Comparative Trust Evaluation}\label{exp:4}
We compare the per-sample PaTAS assessment against three established uncertainty baselines: maximum softmax probability, MC~Dropout~\cite{gal2016dropoutbayesianapproximationrepresenting} ($T{=}30$ passes, $p{=}0.2$), and Evidential Deep Learning~\cite{sensoy2018evidentialdeeplearningquantify}, each trained with the modifications its method requires under the shared architecture, optimizer, and epoch budget; DeepTrust~\cite{10.3389/frai.2020.00054} has no public implementation and is discussed in \cref{sec:rw}. Primary datasets are MNIST and Fashion-MNIST~\cite{xiao2017fashionmnist} (784-512-10, 20 epochs); GTSRB~\cite{GTSRB} (1024-128-43, features standardized with train statistics) serves as an applicability-boundary case discussed below.

The PaTAS assessment instantiates route~(ii), estimation from data, of \cref{sec:acquisition}: each input feature receives a trust opinion from its conformity to the training-data statistics, mapped through the same Baseline-Prior Quantification as parameter trust (\cref{eq:q1}), with per-feature evidence weighted by the feature's reliability as a witness, $w_j = (\sigma_{\mathrm{ref}}/\sigma_j)^2$ (full construction and parameters in Appendix~\ref{app:conformity}). The trust factor is the serial discount $P = P_{\mathrm{model}} \cdot P_{\mathrm{input}}$ of the IPTA path trust under a fully trusted input and the pooled input conformity, and the reported score anchors the discounted confidence at the class prior, $P \cdot \hat{p} + (1-P)/K$. Every method is evaluated identically on clean test data, under Bernoulli-uniform feature corruption ($p \in \{0.3, 0.6\}$, amplitude scaled to the feature range), on a mixed batch pooling clean and corrupted inputs 1:1, and as an OOD detector (the OOD set passes through the identical preprocessing: Fashion-MNIST for MNIST and vice versa); post-hoc isotonic calibration is fit per method on a held-out clean split. Results are mean~$\pm$~std over three evaluation seeds.

\paragraph*{Implementation Details for All Experiments}
Feedforward multiplications are implemented as SL trust discounting and additions as generalized SL averaging fusion (\cref{def:trustfunc}; operational flow in \cref{fig:ptasoperation}); trust revision uses the same fusion. The full Python/NumPy implementation and all experiment scripts are available as \suppp, with instructions for reproducing every experiment.

\vspace{-0.5em}
\subsection{Results and Analysis}
Evaluation is performed during training: after each feedforward/backpropagation iteration we assess the output trustworthiness under three boundary input profiles, fully trusted, fully uncertain, and fully distrusted, tracking the trust, uncertainty, and distrust mass over training, i.e., nine curves per experiment (Experiments 1--3)\footnote{The base rate is set to $0.5$, the maximum-entropy choice over $\{\text{trust}, \text{distrust}\}$, and held constant across all experiments; the score anchoring in \cref{exp:4} uses the class prior $1/K$ where a prior probability of correctness is the natural reference. A systematic sensitivity study of the base rate is left to future work.}. Results are collected in Appendix~\ref{app:expres} and summarized in \cref{tab:cancer_results,tab:mnist,tab:mnistpois,tab:mnistpoisipta}; they confirm several properties proven in \cref{sec:ptastheory}: \emph{stability} (after accuracy converges, the trust assessment settles into a narrow band, the empirical counterpart of the idealized convergence in \cref{thm:convergence}), \emph{vacuity preservation} (a fully uncertain input always produces a fully uncertain output), and \emph{symmetry} (symmetric input trust assessments yield symmetric inference results).

\begin{table*}[ht]
\begin{minipage}{0.63\textwidth}
\centering\footnotesize
\begin{tabular}{cccccc}
\hline
X Trust Assessment & Y Trust Assessment& \(\epsilon = 0.01\) & \(\epsilon = 0.1\) & Train (\%) & Test (\%) \\
\hline
fully distrusted & any & 0 & 0 & 47--58 & 41--55 \\
any & fully distrusted & 0 & 0 & 96--98 & 3--4 \\
fully uncertain & fully uncertain & 0.17 & 0.28 & 66 & 85 \\ % noise-noise
% \hline
fully uncertain & fully trusted & 0.15 & 0.32 & 96 & 96 \\ % noise-clean
% \hline
 fully trusted & fully uncertain & 0.12 & 0.29 & 69 & 94 \\ % clean-noise
% \hline
fully trusted & fully trusted & 0.62 & 0.88 & 98 & 97 \\ % clean-clean
(0.25, 0.25, 0.5) & (0.25, 0.25, 0.5)& 0.06 & 0.11 & 66 & 85 \\
% \hline
(0.25, 0, 0.75) & (0.25, 0, 0.75)& 0.20 & 0.35 & 75 & 94\\
% \hline

\hline
\end{tabular}
\caption{Breast Cancer: final trust mass and train/test accuracy (\%). Rows with a fully distrusted assessment collapse to zero trust and are grouped: distrusted features give train 47--58\% and test 41--55\%, distrusted labels give train 96--98\% and test 3--4\%.}
\label{tab:cancer_results}
\end{minipage}
\hfill
\begin{minipage}{0.37\textwidth}
\centering\footnotesize
\begin{tabular}{cccc}
\hline
Hidden Neurons & Trust \(t\) & Train (\%) & Test (\%) \\
\hline
16 & 0.306 & 65 & 72\\
% \hline
32 & 0.325 & 66 & 89 \\
% \hline
64 & 0.333 & 67 & 90\\
% \hline
128 & 0.337 & 68 & 93\\
% \hline
16 (trust) & 0.778 & 96 & 95\\
\hline
\end{tabular}
\caption{MNIST: final trust mass per architecture.}
\label{tab:mnist}
\end{minipage}

\end{table*}


We therefore report only the fully trusted input case in detail: uncertain inputs yield uncertain outputs, symmetry makes the distrusted case a mirror of the trusted one (\cref{thm:infvac}), and under trusted inputs the distrust mass stays near zero, so the trust mass alone determines the opinion; the tables report its converged value.


\subsubsection*{Experiment 1 -- Breast Cancer Classification}
Corrupted labels mislead rather than prevent learning: training accuracy stays high (96--98\%) while test accuracy collapses to near zero; corrupted features eliminate learning in every setting (41--55\% test, near chance), while label noise depresses training accuracy (69\%) yet still permits generalization (94\%).

We evaluated the system for two values of $\epsilon$ (0.1 and 0.01), the threshold used by the NodeTrust function in \cref{alg:trustupdate}, as summarized in \cref{tab:cancer_results}. For $\epsilon = 0.1$, trust mass under an uncertain $X$ stabilizes around 0.28 for uncertain $Y$ and 0.32 for trusted $Y$; a fully trusted $X$ gives 0.29 for uncertain $Y$ and $0.88$ for trusted $Y$, and whatever the $\epsilon$ value, a distrusted $X$ or $Y$ drives the trust mass rapidly to 0. For $\epsilon = 0.01$ the system reads more gradients as large and all partial-trust configurations drop (0.12--0.17); notably, under degraded features the trusted-$Y$ case falls \emph{below} the uncertain-$Y$ case (0.15 vs.\ 0.17), since trusted labels turn persistent gradient corrections into negative evidence, the mechanism analyzed in \cref{sec:acquisition}.
Overall, trust mass correlates strongly with test accuracy among the non-distrusted configurations. Two findings stand out. First, \emph{label trust matters more than feature trust}: fully uncertain $X$ with trusted $Y$ yields a higher trust mass (0.32) than trusted $X$ with uncertain $Y$ (0.29). Second, \emph{distrust is more damaging than uncertainty}: replacing the fully uncertain opinion $(0,0,1)$ on both $X$ and $Y$ with the mixed assessment $(0.25,0.25,0.5)$ reduces the trust mass from $0.28$ to $0.11$.




\subsubsection*{Experiment 2 -- MNIST Dataset}
PaTAS produces ten output trust opinions, one per class; as they are almost identical we record only the aggregated trust mass. As shown in \cref{tab:mnist}, under noisy features and labels accuracy improves with model size, and test accuracy consistently exceeds training accuracy, suggesting training noise makes the model underestimate its learning progress. Trust also increases with model size, but with diminishing gains (0.306 $\to$ 0.325 $\to$ 0.333 $\to$ 0.337), contrasting sharply with the fully trusted evaluation on the smallest architecture (0.778): \emph{training a smaller architecture on trusted data benefits trust far more than training a larger architecture on uncertain data}. Larger architectures also yield more stable assessments, likely due to smaller average gradient magnitudes.


\subsubsection*{Experiment 3 – Poisoned MNIST Dataset}
PaTAS again produces ten output trust opinions, now informative because the reliability of the inference path varies across labels when some are patched and others are not. \Cref{tab:mnistpois} reports the trust mass for label 3 (clean) and label 6 (poisoned; see \cref{exp:3}), together with clean and patched test accuracy for both digits.

Across all patch sizes the trust mass of label 3 remains higher than that of label 6 by a wide margin (e.g., 0.90 vs.\ 0.58 for the $4\times4$ patch), confirming that PaTAS reliably distinguishes clean from poisoned inference paths. Both values decrease as the patch grows, since larger patches affect more neurons along the inference path, reducing the reliability of the associated parameters and progressively touching even the neurons of clean classes; degradation is largest at $27 \times 27$, where the patch dominates nearly the entire input (0.70 and 0.49), while the clean--poisoned separation remains interpretable at every size.

To evaluate the IPTA, \cref{tab:mnistpoisipta} reports accuracy and trust opinions obtained by feedforwarding a fully trusted opinion through the IPTA derived from clean digit-3, clean digit-6, and patched digit-6 images. The backdoor is near-total in this run (patched 6s are misclassified in 99\% of cases), and the assessment exposes it sharply: clean digit-6 samples retain high accuracy (97.4\%) yet carry a markedly lower trust mass ($0.570$) than clean digit-3 samples ($0.886$), because the parameters along the digit-6 inference path were learned from poisoned data: an accurate prediction whose supporting evidence PaTAS refuses to endorse, precisely the reliability gap that confidence scores do not reveal. Patched samples inherit the same reduced trust and elevated uncertainty ($0.571$, $u = 0.429$), and explicitly distrusting the patch pixels additionally surfaces a non-zero distrust mass ($0.560, 0.011, 0.429$). PaTAS thus provides interpretable warnings about poisoned outputs even when accuracy alone looks healthy.

\begin{table}[ht]
\centering\footnotesize
\setlength{\tabcolsep}{4pt}
\begin{tabular}{ccccccc}
\hline
Patch & \multicolumn{2}{c}{Trust} & \multicolumn{2}{c}{Clean acc.\ (\%)} & \multicolumn{2}{c}{Patched acc.\ (\%)} \\
size & 3 & 6 & 3 & 6 & 3 & 6 \\
\hline
\((1 \times 1)\) & 0.897 & 0.577 & 97.92 & 97.29 & 78.22 & 0.84 \\
\((4 \times 4)\) & 0.895 & 0.576 & 97.82 & 97.39 & 39.60 & 1.04 \\
\((10 \times 10)\) & 0.881 & 0.570 & 97.92 & 97.39 & 22.08 & 0.84 \\
\((27 \times 27)\) & 0.697 & 0.490 & 97.92 & 97.29 & 0 & 0 \\
\hline
\end{tabular}
\caption{Poisoned MNIST; overall train/test accuracy (96.4--99.6\%/97.7--97.8\% across patch sizes) is omitted.}
\label{tab:mnistpois}
\end{table}

\begin{table}[ht]
\centering\footnotesize
\begin{tabular}{ccccccccc}
\hline
 & Accuracy (\%) & Trust  & Distrust & Uncertainty \\
\hline
Clean 3 & 97.82 & 0.886 & 0.0 & 0.114\\

Clean 6  & 97.39 & 0.570 & 0.0 & 0.430\\

6 with patch& 1.04 & 0.571 & 0.0 & 0.429\\

6 with patch & -- & 0.560 & 0.011 & 0.429\\
(patch distrusted) &&&&\\
\hline
\end{tabular}
\caption{IPTA on the poisoned model ($4\times 4$ patch, single run; the backdoor drives patched-6 accuracy to 1\%), feedforwarding a fully trusted opinion; in the last row patch pixels are explicitly distrusted. The poisoned class carries reduced trust even on clean, correctly classified inputs.}
\label{tab:mnistpoisipta}
\end{table}

\subsubsection*{Experiment 4 -- Comparative Trust Evaluation}
\Cref{tab:comparative} reports the two detection tasks in which trust and confidence should differ: recognizing out-of-distribution inputs, and recognizing corrupted inputs inside a mixed batch. PaTAS separates MNIST from Fashion-MNIST at $0.975$ AUROC (FPR@95 $0.145$), against $0.883/0.413$ for softmax and $0.936/0.310$ for MC~Dropout, and detects corrupted inputs at $0.972$ (MNIST) and $0.835$ (Fashion-MNIST) where every confidence-based baseline remains near chance ($0.51$--$0.63$). The ablation row attributes the mechanism: replacing the conformity-based input opinions with the constant fully-trusted opinion collapses PaTAS to baseline level ($0.896/0.415$ OOD, $0.639$ corruption detection), so the input-level trust of \cref{def:trustfunc}, not the confidence signal, carries the detection.

\begin{table}[ht]
\centering\footnotesize
\setlength{\tabcolsep}{3.5pt}
\begin{tabular}{lcccc}
\toprule
 & \multicolumn{2}{c}{OOD: MNIST vs.\ F-MNIST} & \multicolumn{2}{c}{Corruption det.\ $\uparrow$} \\
Method & AUROC $\uparrow$ & FPR@95 $\downarrow$ & MNIST & F-MNIST \\
\midrule
PaTAS (ours) & $\mathbf{0.975}{\scriptstyle\pm.001}$ & $\mathbf{0.145}$ & $\mathbf{0.972}{\scriptstyle\pm.002}$ & $\mathbf{0.835}{\scriptstyle\pm.003}$ \\
\;\;w/o input trust & $0.896$ & $0.415$ & $0.639$ & -- \\
Softmax & $0.883{\scriptstyle\pm.004}$ & $0.413$ & $0.631$ & $0.558$ \\
MC Dropout & $0.936{\scriptstyle\pm.002}$ & $0.310$ & $0.621$ & $0.545$ \\
EDL & $0.883{\scriptstyle\pm.004}$ & $0.658$ & $0.547$ & $0.510$ \\
\bottomrule
\end{tabular}
\caption{Detection of distribution violations (mean $\pm$ std, three seeds): each method's own score separates clean in-distribution inputs from OOD inputs (identical preprocessing), and clean from corrupted inputs in the mixed batch. The ablation row replaces conformity-based input opinions with the constant fully-trusted opinion.}
\label{tab:comparative}
\end{table}

On in-distribution data the two signals coincide: PaTAS matches the softmax filter on correct-vs-incorrect discrimination on clean test data ($0.979 \pm 0.003$ for both, MNIST), with equal calibration error after identical post-hoc calibration ($\mathrm{ECE}=0.006$). Under corruption they separate in the expected direction: accuracy barely drops ($97.9\% \to 97.6\%$ on the mixed batch), so the softmax retains a higher correct-vs-incorrect AUROC ($0.974$ vs.\ $0.895$) precisely because it ignores the distribution violation that PaTAS flags. Trust and confidence are thus complementary: confidence ranks correctness \emph{within} the training distribution, while trust exposes inputs that \emph{leave} it, where confident predictions are unwarranted even when correct. Directionality follows: in the reverse OOD task (Fashion-MNIST in-distribution, MNIST as OOD), MNIST digits occupy a quiet subset of Fashion-MNIST's marginals rather than violating them, and PaTAS ($0.781$) offers no advantage over MC~Dropout ($0.834$); conformity-based trust detects support \emph{violations}, not support \emph{subsets}.

\begin{table}[ht]
\centering\footnotesize
\setlength{\tabcolsep}{4.5pt}
\begin{tabular}{lccccc}
\toprule
Dataset & Med.\ $\sigma_{\mathrm{eff}}$ & $w_{\max}$ & Top-10\% & $w{\geq}4$ & Corr.\ det. \\
\midrule
MNIST & 0.47 & 52.3 & 0.32 & 0.40 & 0.972 \\
Fashion-MNIST & 0.90 & 64.0 & 0.76 & 0.18 & 0.835 \\
GTSRB & 0.98 & 1.8 & 0.14 & 0.00 & 0.545 \\
CIFAR-10 (gray) & 0.23 & 1.1 & 0.11 & 0.00 & -- \\
\bottomrule
\end{tabular}
\caption{A-priori applicability diagnostics: conformity-model statistics on each dataset's clean training features \emph{before any model is trained} (quantities defined in Appendix~\ref{app:conformity}), against the corruption detection later observed. Registered domains concentrate evidence in reliable witnesses; unregistered imagery has none, so input-level conformity is predicted, and for GTSRB confirmed, to be uninformative.}
\label{tab:diagnostics}
\end{table}

\Cref{tab:diagnostics} explains when this input-level signal exists. Fitting the conformity model on each dataset's training features yields a self-diagnostic: spatially registered domains (MNIST, Fashion-MNIST) concentrate evidence in a minority of tight-marginal ``reliable witness'' features, whose fraction even predicts the ordering of the observed detection margins ($0.40 \to 0.972$, $0.18 \to 0.835$, $0.00 \to$ chance). Unregistered imagery (GTSRB, CIFAR-10) has uniformly wide marginals, so no feature is a reliable witness and marginal conformity is predicted to be uninformative, which the full GTSRB battery confirms (corruption detection $0.545$; OOD vs.\ grayscale CIFAR-10 $0.661$ vs.\ $0.669$ for softmax). Applicability can thus be established per dataset before deployment; richer input-trust models for unregistered domains are future work.

Finally, the model-side trust provides a training-quality audit that requires \emph{no test data at all}. Training the MNIST network under increasing label-flip rates $p$ and reading the PaTAS feedforward trust under a fully trusted input (\cref{tab:audit}) shows a tripwire behavior: trust drops by $31.7\%$ relative to the clean-trained model at \emph{any} non-zero corruption rate, including $p{=}0.1$ where test accuracy moves by only $3$ points and mean confidence by $14\%$. The flat level across rates is the designed saturation of the evidence counting: already at $p{=}0.1$, a 128-sample mini-batch contains corrupted labels with probability $1-0.9^{128} \approx 1-10^{-6}$, so every batch delivers the persistent large-gradient corrections that \textsc{NodeTrust} counts as negative evidence, and the converged trust registers \emph{that} corruption is present rather than how much. Mean softmax confidence instead tracks the flip rate almost linearly and \emph{measures} severity, but reading it requires clean held-out inputs (and labels, if it is to be compared against accuracy), and a moderate value is indistinguishable from a merely difficult task. The two signals answer different questions: PaTAS detects \emph{that} training data was corrupted from the training dynamics alone; confidence quantifies the damage given labelled data.

\begin{table}[ht]
\centering\footnotesize
\setlength{\tabcolsep}{5pt}
\begin{tabular}{cccccc}
\toprule
$p$ & Test acc.\ (\%) & Mean conf. & PaTAS trust & $\Delta$trust & $\Delta$conf. \\
\midrule
0   & 97.9 & 0.981 & 0.998 & --     & --     \\
0.1 & 94.7 & 0.841 & 0.682 & 31.7\% & 14.2\% \\
0.2 & 95.7 & 0.762 & 0.682 & 31.7\% & 22.3\% \\
0.3 & 93.7 & 0.684 & 0.682 & 31.7\% & 30.3\% \\
0.4 & 91.2 & 0.589 & 0.682 & 31.7\% & 39.9\% \\
0.5 & 82.8 & 0.476 & 0.681 & 31.7\% & 51.5\% \\
\bottomrule
\end{tabular}
\caption{Training-quality audit (MNIST, 784-512-10) under label-flip rate $p$. PaTAS trust (the feedforward opinion under a fully trusted input, no test data needed) acts as a corruption tripwire ($\Delta$ = relative drop vs.\ $p{=}0$); mean softmax confidence tracks severity but requires clean held-out data.}
\label{tab:audit}
\end{table}
